\section{Compact spatial equations through the third threshold}
\label{tgs:section}

The third growth starts at the actual second return $t_b=t_2^b$.
The preceding coefficient construction proves its actual crossing time
$t_{a3}:=t_3^a>t_b$ in Theorem~\ref{tge:actual_threshold} of the
original threshold. This section realizes that
construction in the full equations, with the same physical diffusivity
$d>0$ in both equations and the declared pressure $p=0$:
\[
\begin{aligned}
 \theta_t+u\cdot\nabla\theta&=d\Delta\theta+f_\theta,\\
 u_t+(u\cdot\nabla)u&=d\Delta u+\theta e_2+f_u,
 &\operatorname{div}u&=0.
\end{aligned}
\]
Every field is defined on $\mathbb R^2$ in the original coordinate $x$.
Retain
\[
 J=\begin{pmatrix}0&-1\\1&0\end{pmatrix},\quad
 A_0\geq1,\quad\nu=\sqrt{A_0},\quad
 \mu=\lambda_0=\lceil\sqrt{A_0}/2\rceil,\quad
 K_j=\mu\widehat\lambda_j\quad(1\leq j\leq3).
\]
Here $\nu$ is the source time scale, not $d$. The hats distinguish
the source frequencies from their physical multiples $K_j$.
The released recursive frequency, source growth logarithm, inherited
scales, angle, and activation are
\begin{equation}\label{tgs:scales}
 \begin{gathered}
 \widehat\lambda_3=\widehat\lambda_2^{Q_3},\qquad
 L_3=(k_3+5+1/8)\log\widehat\lambda_3,\qquad
 s_3=L_3\sigma_1/\sigma_2,\qquad\Gamma_3=\sigma_2\sin s_3,\\
 S_3=(\nu^2/\mu)\widehat\lambda_3^{-k_3-6},\qquad
 P_3^*=S_3e^{L_3}=(\nu^2/\mu)\widehat\lambda_3^{-7/8},\qquad
 w_3(t)=\eta(\Gamma_3(t-t_b)).
 \end{gathered}
\end{equation}
The source step $\eta$ is smooth, nondecreasing, zero on
$(-\infty,0]$, and one on $[1,\infty)$.
The actual coefficient proof, rather than an inviscid comparison,
supplies the following bounds through $t_{a3}$ and existence through
the larger horizon $t_b+16/\sigma_1$:
\begin{equation}\label{tgs:actualbounds}
 \begin{gathered}
 h_0\leq h<3,\quad 0<\Omega_1<9\sigma_1,\quad
 1/4<B\leq2,\quad a=\sin s_2\in(0,1/8],\quad
 h_0=\cos s_2\geq\sqrt{63}/8,\\
 \chi=(h^2+a^2)^{1/2},\quad r_1=1,\quad
 1\leq\chi^2<10,\quad 1\leq r_3^2\leq2,\quad
 0< -\Theta_3\leq P_3^*,\quad
 |\Omega_3|\leq3K_3P_3^*/\sigma_2.
 \end{gathered}
\end{equation}
The angle branch is the actual small source branch $0<s_2,s_3<1/8$.
The old amplitudes remain the actual ones:
\[
 \Theta_1=-\sigma_1^2B/K_1,\qquad
 \Omega_2=0,\qquad
 \Theta_2(t)=-P_{2b}\exp\!\left[-dK_2^2
                 \int_{t_b}^{t}\chi(r)^2\,dr\right],
 \qquad0<P_{2b}\leq e^{18}P_2^*.
\]
These equations do not hold either old temperature fixed.

\subsection{The exact third material flow and support inclusions}

Put $\vartheta=\arctan(a/h)$ and write $h_b=h(t_b)$,
$\vartheta_b=\vartheta(t_b)$. The actual held second feedback is
\begin{equation}\label{tgs:parents}
 \begin{gathered}
 \alpha=s_*+\vartheta,\qquad
 \dot h=a\Omega_1,\qquad
 \dot\alpha=-a^2\Omega_1/\chi^2,\\
 \zeta_1=(-a,h)/\chi,\qquad\zeta_2=(0,\chi),\\
 G_{<3}=-A_0R_\alpha e_2+K_1\Theta_1\zeta_1
                                  +K_2\Theta_2\zeta_2,\\
 D_{<3}=\dot\alpha J+\Omega_1J\zeta_1\otimes\zeta_1.
 \end{gathered}
\end{equation}
The second velocity increment is zero, but its scalar increment
contributes the entire last term in $G_{<3}$.
The closed older amplitude equations on this interval are
\begin{equation}\label{tgs:olderpair}
 \begin{aligned}
 \dot\Theta_1&=\frac{A_0\sin s_*}{K_1}\Omega_1-dK_1^2\Theta_1,
 &\dot\Omega_1&=-\frac{aK_1}{\chi}\Theta_1-dK_1^2\Omega_1,\\
 \dot\Theta_2&=-dK_2^2\chi^2\Theta_2,
 &\dot\Omega_2&=0.
 \end{aligned}
\end{equation}
Indeed $-K_1\Theta_1=A_1$ in \eqref{tge:older_hold}
gives the first equation, and $\zeta_{1,1}=-a/\chi$ in
\eqref{ms:amplitudes} gives the second. The held identities
$\Omega_2=\zeta_{2,1}=0$ give the remaining two equations.
For $S(q)=\left(\begin{smallmatrix}1&-q\\0&1\end{smallmatrix}\right)$,
the matrix in the exact geometry calculation is
\begin{equation}\label{tgs:flow}
 M_3(t)=R_\vartheta S((h-h_b)/a)R_{-\vartheta_b}
       =M_2(t)M_2(t_b)^{-1},\qquad
 \dot M_3=D_{<3}M_3,\qquad M_3(t_b)=I.
\end{equation}
Its determinant is one. Its inverse is
$R_{\vartheta_b}S(-(h-h_b)/a)R_{-\vartheta}$.
Indeed the rotation matrices have inverse rotations and
$S(q)^{-1}=S(-q)$, while differentiating $S$ gives
$\dot S=\Omega_1(J e_2\otimes e_2)S$.
Since $\zeta_1=R_\vartheta e_2$, multiplication proves the ODE
in \eqref{tgs:flow} directly, including its signs.
Integrating $\dot h=a\Omega_1$ over $16/\sigma_1$ gives
$0\leq(h-h_b)/a\leq144$. The elementary inequality
$\|S(q)\|\leq1+|q|$, obtained from
$S(q)=I-qe_1\otimes e_2$, proves
\begin{equation}\label{tgs:matrixbounds}
 \|M_3\|,\|M_3^{-1}\|\leq N_3:=145.
\end{equation}
Use the original unit insertion vector $p_3=e(s_3)$ and put
$\zeta_3=M_3^{-T}p_3$. Differentiating this exact inverse proves
$\dot\zeta_3=-D_{<3}^T\zeta_3$. Thus the moving phase and
moving material envelope below are transported by the same actual
older affine velocity, without changing the insertion coordinates.

The released functions $F,P,f$ are retained without alteration:
$F(s)=\sin s+\chi_F(s)(s-\sin s)$ on $[-\pi,\pi]$ and
periodically extended, with $\chi_F$ smooth and even,
$F(s)=s$ on $|s|\leq s_0$, $0<s_0<1$;
$P'=F$, $P$ is even and periodic, and
$\int_{-\pi}^{\pi}P=0$.
The radial $f\in C_c^\infty(B_1)$ satisfies $0\leq f\leq1$
and $f=1$ on $B_{1/2}$. In particular the constant $P(0)$
on the affine cell is preserved. Keep the source radii
\begin{equation}\label{tgs:radii}
 \ell_1=\frac{s_0}{C_\ell\mu},\quad
 \ell_2=\frac{\widehat\lambda_1^{-3}}\mu,\quad
 \ell_3=\frac{\widehat\lambda_2^{-3}}\mu,\qquad C_\ell\geq4.
\end{equation}
Let $N=1+(3-h_0)/a$, so $\|M_2^{\pm1}\|\leq N$.
The actual source frequency selection in \eqref{tge:Y3} proves
the strict inequalities established in the
\hyperref[tge:nesting]{exact third support nesting calculation}:
\begin{equation}\label{tgs:nesting}
 \begin{gathered}
 N\widehat\lambda_1^{-3}<s_0/(2C_\ell),\qquad
 N\widehat\lambda_1^{-2}<s_0,\\
 2N_3N\widehat\lambda_2^{-3}<\widehat\lambda_1^{-3},
 \qquad N_3\sqrt{10}\widehat\lambda_2^{-2}<s_0.
 \end{gathered}
\end{equation}
Here they have an exact spatial meaning. Put
\[
 y_3=M_3^{-1}x,\qquad g_3=f(y_3/\ell_3),\qquad
 k_3^{\rm ph}=K_3\zeta_3,\qquad s_3^{\rm ph}=k_3^{\rm ph}\cdot x.
\]
The superscript distinguishes a phase vector and a phase variable
from the source seed exponent $k_3$ and insertion angle $s_3$.
On the support of any third increment or derivative,
$|x|\leq N_3\ell_3$. Consequently
\[
 |M_2^{-1}x|\leq NN_3\ell_3<\ell_2/2,\qquad
 |K_2\zeta_2\cdot x|\leq K_2\sqrt{10}N_3\ell_3<s_0.
\]
Thus the second envelope is one and its profile is exactly affine.
Moreover $x=M_2y_2$ with $|y_2|<\ell_2/2$, whence
$|x|<N\ell_2/2$. The first two inequalities in
\eqref{tgs:nesting} imply $|x|<\ell_1/2$ and
$|K_1\zeta_1\cdot x|<s_0$.
This proves inclusion in the first affine cell as well.
The base is affine there because $\mu\ell_1=s_0/C_\ell<s_0$.
All inequalities are strict, so the identities hold on an open
neighborhood of the supports; differentiating a force introduces
no missing boundary term. The actual older fields are therefore
exactly $G_{<3}\cdot x$ and $D_{<3}x$ at every point needed below.

\subsection{Third fields and the entire scalar and vector forces}

Write $k=k_3^{\rm ph}$, $\rho=|k|^2=K_3^2r_3^2$, $g=g_3$,
$s=s_3^{\rm ph}$, $D=D_{<3}$, and $G=G_{<3}$ in this
subsection only. The actual third pair satisfies
\begin{equation}\label{tgs:pair}
 \dot\Theta_3=-\frac{Jk\cdot G}{\rho}\Omega_3-d\rho\Theta_3,
 \qquad\dot\Omega_3=k_1\Theta_3-d\rho\Omega_3.
\end{equation}
Define, with all one-variable profiles evaluated at $s$,
\begin{equation}\label{tgs:fields}
 T=w_3\Theta_3,\quad Q=\frac{w_3\Omega_3}{\rho},\quad
 V_3=TFg,\quad\Psi_3=QPg,\quad
 v_3=J\nabla\Psi_3=QJ(kFg+P\nabla g).
\end{equation}
For $t\leq t_b$ these three increments are defined to be zero.
The flat activation and the smooth finite ODE solution make this
zero extension smooth in space and time. On $[0,t_{a3}]$ set
\[
 \theta=\theta_b+V_1+V_2+V_3,\qquad
 u=u_b+v_1+v_2+v_3,\qquad p=0.
\]
The base and first two increments are precisely those in
\eqref{ms:base}--\eqref{ms:fields}, continued by their actual
holding equations through the third growth. The initial data
remain exactly
$\theta(x,0)=-(A_0/\mu)\sin(\mu x_2)\chi_0(x)$ and $u(x,0)=0$.
Define
\begin{equation}\label{tgs:Qrest}
 Q_* =\frac{\dot w_3\Omega_3+w_3k_1\Theta_3}{\rho}
                         -Q\frac{\dot\rho}{\rho},\qquad
 \dot\rho=-2k\cdot D^Tk.
\end{equation}
Then $\dot Q=Q_*-d\rho Q$, including the varying phase length.
The complete third scalar and full vector force increments are
\begin{align}
 E_{\theta,3}={}&\dot w_3\Theta_3Fg+QP(J\nabla g\cdot G)
       +QT(F^2-PF')g(Jk\cdot\nabla g)\nonumber\\
 &-dT\{\rho(F+F'')g+2F'k\cdot\nabla g+F\Delta g\},
 \label{tgs:scalar}\\
 E_{u,3}={}&2Dv_3+Q_*J(kFg+P\nabla g)
        +(\nabla v_3)v_3-TFg e_2-d(\rho v_3+\Delta v_3).
 \label{tgs:vector}
\end{align}
In components, every derivative required in these formulas is
\begin{align}
 \nabla v_3={}&QJ\{k\otimes k F'g
       +F(k\otimes\nabla g+\nabla g\otimes k)+P\nabla^2g\},
 \label{tgs:velocityjet}\\
 \Delta V_3={}&T\{\rho F''g+2F'k\cdot\nabla g+F\Delta g\},
 \label{tgs:scalarLaplacian}\\
 \Delta v_3={}&QJ\{\rho kF''g+\rho F'\nabla g
       +2F'k(k\cdot\nabla g)+2F(\nabla^2g)k\nonumber\\
 &\hspace{38mm}+Fk\Delta g+P\nabla\Delta g\}.
 \label{tgs:vectorLaplacian}
\end{align}
The total forces are
\begin{equation}\label{tgs:totalforce}
 f_\theta=f_{\theta,b}+\sum_{j=1}^3E_{\theta,j},\qquad
 f_u=f_{u,b}+\sum_{j=1}^3E_{u,j},
\end{equation}
with the exact base and old increments of
\eqref{ms:scalar}--\eqref{ms:baseforce} evaluated at the actual
continued coefficients. There is no change of pressure hidden in
the vector formula.

\begin{proof}[Proof of both equations and every interaction direction]
For $\mathcal D=\partial_t+Dx\cdot\nabla$,
\eqref{tgs:flow} and its inverse give
$\mathcal D(M_3^{-1}x)=0$, and the phase equation gives
$\mathcal D(k\cdot x)=0$. Hence $\mathcal Dg=0$,
$\mathcal DV_3=\dot T Fg$, and $\mathcal D\Psi_3=\dot QPg$.
Subtracting the old scalar equation from the new one gives
$\dot T Fg+v_3\cdot G+v_3\cdot\nabla V_3-d\Delta V_3$.
The two exact contractions are
\[
 v_3\cdot G=QFg(Jk\cdot G)+QP(J\nabla g\cdot G),\qquad
 v_3\cdot\nabla V_3=QT(F^2-PF')g(Jk\cdot\nabla g).
\]
Substitution of \eqref{tgs:pair} cancels the first contraction
against the coupling in $\dot T$ and proves \eqref{tgs:scalar}.

Since $\operatorname{tr}D=0$, direct multiplication gives
$DJ+JD^T=0$. Therefore
$\mathcal D(J\nabla\Psi_3)=Dv_3+J\nabla(\mathcal D\Psi_3)$.
The reverse interaction is $(v_3\cdot\nabla)(Dx)=Dv_3$.
Together with self-advection, buoyancy, and diffusion, this proves
\eqref{tgs:vector} after substituting \eqref{tgs:Qrest}.
The ordinary product rule proves
\eqref{tgs:velocityjet}--\eqref{tgs:vectorLaplacian}; in particular
the last display contains all phase, envelope, and mixed terms.

The operator $\mathcal D$ uses the full old velocity $u_b+v_1+v_2$,
which equals $Dx$ on the third support. It accounts for
$(u_b+v_1+v_2)\cdot\nabla V_3$ and
$(u_b+v_1+v_2)\cdot\nabla v_3$.
The contractions with $G$ and $D$ account for
$v_3\cdot\nabla(\theta_b+V_1+V_2)$ and
$v_3\cdot\nabla(u_b+v_1+v_2)$.
Although $v_2=0$, $v_3\cdot\nabla V_2$ generally is not zero;
it is exactly the $K_2\Theta_2\zeta_2$ part of $v_3\cdot G$.
The earlier increments already account for interactions among
the base and first two layers. Outside a third support all its
derivatives vanish. Adding these identities therefore proves
\eqref{tgs:totalforce} globally, with every ordered pair counted
once. Finally each velocity is a rotated gradient, so mixed
partial symmetry proves $\operatorname{div}u=0$.
\end{proof}

For the exact coordinate correspondence put
$C_3=M_3^{-1}M_3^{-T}$ and $\widetilde g(y)=f(y/\ell_3)$.
The spatial gradient and Laplacian pulled back to $(s,y)$ are
\begin{equation}\label{tgs:metric}
 \begin{split}
 \mathcal P_3&=K_3\zeta_3\partial_s+M_3^{-T}\nabla_y,\\
 \mathcal L_3&=K_3^2r_3^2\partial_s^2
       +2K_3(C_3p_3\cdot\nabla_y)\partial_s
       +\sum_{l,m}(C_3)_{lm}\partial_{y_l}\partial_{y_m}.
 \end{split}
\end{equation}
The chain rule gives the first identity, and multiplying its
commuting spatial derivatives gives the second. In particular
$k\cdot\nabla_xg=K_3(C_3p_3)\cdot\nabla_y\widetilde g$ and
$\Delta_xg=\operatorname{tr}(C_3\nabla_y^2\widetilde g)$.
This proves the original-coordinate meaning of all mixed terms,
without imposing a scalar metric or a fixed cutoff.

\subsection{All spatial derivative norms and diffusion costs}

For any scalar, vector, or tensor $U$, set
$|U|_m=\sup_x\|D_x^mU(x)\|$ in the Euclidean multilinear
operator norm, $m\geq0$. For one-variable $W$ set
$[W]_m=\sup_s|W^{(m)}(s)|$, and set
$f_m=\sup_y\|D^mf(y)\|$.
Define only explicit, finite constants
\begin{align}
 \kappa_3^*&=\sqrt2K_3,&
 e_{3,m}&=(N_3/\ell_3)^mf_m,\nonumber\\
 \mathcal B_{3,m}(W)&=\sum_{q=0}^m\binom mq
                (\kappa_3^*)^q[W]_qe_{3,m-q},&
 \mathcal C_{3,m}(P)&=\sum_{q=0}^m\binom mq
                (\kappa_3^*)^q[P]_qe_{3,m-q+1}.
 \label{tgs:normdefinitions}
\end{align}
The symbol $\kappa_3^*$ is a frequency bound, not a diffusivity.
Repeated product and chain rules prove
\[
 |g_3|_m\leq e_{3,m},\qquad
 |W(s_3^{\rm ph})g_3|_m\leq\mathcal B_{3,m}(W),\qquad
 |P(s_3^{\rm ph})\nabla g_3|_m\leq\mathcal C_{3,m}(P).
\]
Indeed choose the $q$ of $m$ directional derivatives assigned to
the profile, producing $\binom mq$ terms, each bounded by the
corresponding summand. The derivative of a linear phase is $k$,
and that of $M_3^{-1}x/\ell_3$ is $M_3^{-1}/\ell_3$;
their higher spatial derivatives vanish. This proves the stated
operator norm bounds for arbitrary directions.

The actual coefficient bounds give
\begin{equation}\label{tgs:coefficientbounds}
 \begin{gathered}
 U_\Theta=P_3^*,\qquad U_\Omega=3K_3P_3^*/\sigma_2,
 \qquad W_1=\Gamma_3[\eta']_0,\\
 D^*=9\sigma_1(1+a^2),\qquad
 G^*=A_0+2\sigma_1^2+\sqrt{10}K_2e^{18}P_2^*,\\
 Q^*=\frac{U_\Omega}{K_3^2},\qquad
 Q_*^*=\frac{W_1U_\Omega+\kappa_3^*U_\Theta+2D^*U_\Omega}{K_3^2}.
 \end{gathered}
\end{equation}
In fact $\|D\|\leq|\dot\alpha|+\Omega_1$ and
$|\dot\alpha|=a^2\Omega_1/\chi^2\leq9a^2\sigma_1$,
which prove the bound for $D$. The three summands in $G$
give its bound. The denominator $\rho\geq K_3^2$ and
$|\dot\rho/\rho|\leq2D^*$ prove the bounds for $Q,Q_*$.
In particular activation contributes its exact physical factor
$\Gamma_3[\eta']_0$ to both forces.

The full damping-plus-diffusion costs are bounded by
\begin{align}
 \mathcal D_{\theta,3,m}
 &=dU_\Theta\{(\kappa_3^*)^2\mathcal B_{3,m}(F)
                              +2\mathcal B_{3,m+2}(F)\},\nonumber\\
 \mathcal D_{u,3,m}
 &=dQ^*\{(\kappa_3^*)^2\mathcal B_{3,m+1}(P)
                              +2\mathcal B_{3,m+3}(P)\}.
 \label{tgs:diffusioncosts}
\end{align}
These include $-d\rho V_3$ and $-d\rho v_3$ from modal
damping as well as the Laplacians of the changed fields.
The factor two bounds the two coordinate traces defining
$\Delta$ in dimension two. Keeping this factor explicitly avoids
identifying a Laplacian with a single second derivative.
The complete third force estimates are
\begin{align}
 |E_{\theta,3}|_m\leq{}&W_1U_\Theta\mathcal B_{3,m}(F)
           +Q^*G^*\mathcal C_{3,m}(P)+\mathcal D_{\theta,3,m}
 \nonumber\\
 &+Q^*U_\Theta\sum_{q=0}^m\binom mq
       \mathcal B_{3,q+1}(P)\mathcal B_{3,m-q+1}(F),
 \label{tgs:scalarcost}\\
 |E_{u,3}|_m\leq{}&(2D^*Q^*+Q_*^*)\mathcal B_{3,m+1}(P)
           +U_\Theta\mathcal B_{3,m}(F)+\mathcal D_{u,3,m}
 \nonumber\\
 &+(Q^*)^2\sum_{q=0}^m\binom mq
       \mathcal B_{3,q+1}(P)\mathcal B_{3,m-q+2}(P).
 \label{tgs:vectorcost}
\end{align}
The nonlinear scalar term is bounded as $v_3\cdot\nabla V_3$,
and the nonlinear vector term as $(\nabla v_3)v_3$.
The product rule with the displayed allocation of derivatives
proves the two sums. All other terms follow individually from
\eqref{tgs:scalar}--\eqref{tgs:vectorLaplacian} and
\eqref{tgs:normdefinitions}; hence the estimates include every
activation, parent-interaction, self-interaction, buoyancy and
diffusion term in the full vector equation.
For time integrals one has, in particular,
\begin{equation}\label{tgs:integratedcost}
 \int_{t_b}^{t_{a3}}|\text{third damping-plus-diffusion force}|_m\,dt
 \leq\frac{16}{\sigma_1}\mathcal D_{\bullet,3,m},
\end{equation}
with the matching scalar or vector bound. The identical physical
length factor integrates either full-force estimate.

All frequency and localization dependences are retained:
\begin{equation}\label{tgs:frequencycost}
 \frac{N_3}{\ell_3}=145\mu\widehat\lambda_2^3,
 \quad\kappa_3^*=\sqrt2\mu\widehat\lambda_3,
 \quad dU_\Theta K_3^2=d\nu^2\mu\widehat\lambda_3^{9/8},
 \quad dQ^*K_3^3=\frac{3d\nu^2\mu}{\sigma_2}
                      \widehat\lambda_3^{9/8}.
\end{equation}
The last two equalities identify coefficients occurring in the
costs, not lower bounds for the entire forces. No cancellation of
unrelated vector terms has been assumed. The finite positive-$d$
range from the actual coefficient theorem is not replaced here by
a bound uniform over an infinite frequency sequence.

For completeness, the bounds for the old and base forces remain
fully quantitative on the slab ending at the third threshold. The formulas
\eqref{ms:Bnorm}--\eqref{ms:fullvectornorm} hold with the original
$N_1=1,N_2=N$, radii, and global amplitude bounds
\eqref{ms:amplitudebounds}; on the added interval the sharper
$|\Theta_1|\leq2\sigma_1^2/K_1$, $|\Omega_1|\leq9\sigma_1$,
$|\Theta_2|\leq e^{18}P_2^*$, $\Omega_2=0$ prove their continued
validity. If $A_1^{\rm old},A_2^{\rm old}$ are the explicit earlier
control derivative bounds, one may use the following bounds on this slab:
\begin{equation}\label{tgs:oldcontrolbounds}
 \begin{split}
 A_1^{\rm ext}&=\max\{A_1^{\rm old},9a^2\sigma_1\},\\
 A_2^{\rm ext}&=\max\{A_2^{\rm old},
              a^2(4\sigma_1^2+9dK_1^2\sigma_1)
                         +486a^3\sigma_1^2\}.
 \end{split}
\end{equation}
To verify the second number, the holding equations give
$\dot\Omega_1=\sigma_1^2aB/\chi-dK_1^2\Omega_1$,
so $|\dot\Omega_1|\leq4\sigma_1^2+9dK_1^2\sigma_1$.
Differentiating $\dot\alpha=-a^2\Omega_1/\chi^2$ and using
$\dot\chi^2=2ha\Omega_1$, $h<3$, and $\chi^2\geq1$
gives the stated additional term $486a^3\sigma_1^2$.
Substitution into the explicit base norms \eqref{ms:basenorm}
and old force norms completes the estimates for both total forces
in \eqref{tgs:totalforce}, at every spatial order. Their integrals
use this slab's actual interval length $t_{a3}$.

\subsection{Mixed derivatives, affine core, and the threshold state}

No stationary-envelope convention is used to estimate time
derivatives. Set $A(t)=M_3(t)^{-1}/\ell_3$, introduce independent
variables $(t,x,s,y)$, and define
\begin{equation}\label{tgs:jetoperators}
 \mathfrak T=\partial_t+(\dot k\cdot x)\partial_s
                       +(\dot A x)\cdot\nabla_y,\qquad
 \mathfrak X_i=\partial_{x_i}+k_i\partial_s
                       +\sum_lA_{li}\partial_{y_l}.
\end{equation}
For $\iota^*R=R(t,x,k(t)\cdot x,A(t)x)$, the chain rule gives
\begin{equation}\label{tgs:jetidentity}
 \partial_t^b\partial_x^\gamma\iota^*R
   =\iota^*(\mathfrak T^b\mathfrak X_1^{\gamma_1}
                              \mathfrak X_2^{\gamma_2}R).
\end{equation}
Induction on $b+|\gamma|$ proves this equality. The commutators
vanish: $\dot k_i$ cancels the $x_i$ derivative of
$\dot k\cdot x$, and $\dot A_{li}$ cancels the corresponding
derivative of $\dot A x$. This also proves consistency of the
order of mixed derivatives.
Apply these operators to the exact third fields and forces,
starting for example with $Q(t)P(s)f(y)$ for the streamfunction.
Every resulting term is a finite product of coefficient jets,
explicit powers $x^\beta$, profile derivatives, and envelope
derivatives. Its absolute value is bounded by replacing those
factors by their coefficient suprema on $[t_b,t_{a3}]$,
$(N_3\ell_3)^{|\beta|}$, $[P]_q,[F]_q$, and $f_m$,
then summing all terms produced by \eqref{tgs:jetidentity}.
The coefficient jets themselves are evaluated recursively by
\eqref{tgs:parents}, \eqref{tgs:olderpair},
\eqref{tgs:flow}, \eqref{tgs:pair}, and
the explicit activation. Product and quotient rules, in increasing
time order, use only the positive denominators
$a,K_3,\chi^2,\rho$ and the fixed source scales.
Their lower bounds in \eqref{tgs:actualbounds} therefore yield
a finite explicit bound at every prescribed mixed order. The same
recursion from \eqref{ms:jetops}--\eqref{ms:jetpullback} handles
the base and old layers with the extended control bounds. This
procedure includes every time derivative of the diffusion terms
and of the changing material metric.

On a third affine cell, where $g_3=1$ and $|s_3^{\rm ph}|<s_0$,
the retained primitive is exactly $P(s)=P(0)+s^2/2$. Thus
$\Psi_3=Q(P(0)+s^2/2)$ there; taking its gradient, with the
original constant still present in the full cutoff formula, gives
\begin{equation}\label{tgs:affinecell}
 V_3=T K_3\zeta_3\cdot x,\qquad
 v_3=\frac{w_3\Omega_3}{r_3^2}
                  J\zeta_3\otimes\zeta_3\,x,
 \qquad\Delta V_3=\Delta v_3=0.
\end{equation}
The damping parts of the forces are still
$-dK_3^2r_3^2 V_3$ and $-dK_3^2r_3^2 v_3$ there.
On an uncut sine-profile region instead, $F''=-F$ makes these
damping-plus-diffusion contributions zero exactly. The remaining
vector terms in \eqref{tgs:vector} are retained at $p=0$.
This proves precisely how the actual damped pair maps into the
released profile; it is not an unforced eigenfunction assertion.

At the actual third threshold,
\[
 \Theta_3(t_{a3})=-P_3^*,\qquad \Omega_3(t_{a3})<0,
\]
with the exact phase $M_3(t_{a3})^{-T}e(s_3)$, the exact moving
cutoff, and $w_3(t_{a3})=1$. The last equality follows from
\[
 \Gamma_3(t_{a3}-t_b)>L_3/3>1,
\]
proved in
Theorem~\ref{tge:actual_threshold}.
In the common affine cell the complete fields are
\begin{equation}\label{tgs:thresholdfields}
 \theta=(G_{<3}+w_3K_3\Theta_3\zeta_3)\cdot x,
 \qquad
 u=\left(D_{<3}+\frac{w_3\Omega_3}{r_3^2}
                       J\zeta_3\otimes\zeta_3\right)x.
\end{equation}
These are actual states of the full PDE just proved. The third
vorticity has not returned to zero at this threshold; no third
steering or return statement follows from the growth calculation.

\subsection{Smooth finite-slab realization and honest force continuation}

Let $R_{\rm base}$ denote the original spatial support radius
called $R_H$ in \eqref{ms:base}, as distinct from the temporal
horizon $16/\sigma_1$ of the third-growth construction.
The fields and forces are smooth and spatially supported in the
same $B_{R_{\rm base}}$ as the base. All third derivatives are supported
within $|x|\leq N_3\ell_3$; the proved inclusions put this ball
inside the old base support. The velocity has finite energy since
it is smooth and compact. For
$N_0(x)=(2\pi)^{-1}\log|x|$, the compact smooth streamfunction
$\psi=\psi_b+\Psi_1+\Psi_2+\Psi_3$ satisfies
$N_0*\Delta\psi=(\Delta N_0)*\psi=\psi$ by distributional
integration by parts. Applying $J\nabla$ gives exactly
$u=J\nabla N_0*\omega$, $\omega=\Delta\psi$.
Odd $F$, even $P$, and even material envelopes give odd
$\theta,u,f_\theta,f_u$ and even $\omega$.

There is a smooth right collar at the threshold. The actual
coefficient theorem continues the same $F_2$ growth equations
through $t_b+16/\sigma_1$, with $t_{a3}$ strictly before that
horizon. By \eqref{tge:times} its remaining interval exceeds
$2/\sigma_1$. Choose the explicit positive collar length
$\epsilon_R=1/\sigma_1$.
Continue the exact coefficient solution and fields to
$t_{a3}+\epsilon_R$; do not freeze their nonstationary endpoint.
The upper rate in \eqref{tge:gain_rate} holds on that entire
horizon, as does the quotient bound \eqref{tge:quotient}.
Consequently on the collar one has the explicit bounds
\begin{equation}\label{tgs:collaramplitudes}
 -\Theta_3\leq e^{3\Gamma_3/\sigma_1}P_3^*,\qquad
 |\Omega_3|\leq
       \frac{3K_3}{\sigma_2}e^{3\Gamma_3/\sigma_1}P_3^*.
\end{equation}
Indeed integrate $(\log(-\Theta_3))'<3\Gamma_3$ from the
threshold to each collar time and then use the quotient bound.
All spatial collar costs therefore follow from
\eqref{tgs:coefficientbounds}--\eqref{tgs:vectorcost} by replacing
$U_\Theta,U_\Omega$ by the right sides of
\eqref{tgs:collaramplitudes}; the same coefficient-jet recursion
proves the mixed costs. This does not impose the pre-threshold
bound $-\Theta_3\leq P_3^*$ after the crossing.
The strict support inequalities continue on the proved horizon.
Extend the original initial state constantly to the left of zero,
with zero angle and inactive increments as before, and put
$\epsilon_L=1/\nu$.
Multiplying the extended total residual forces by
\begin{equation}\label{tgs:timecutoff}
 \beta(t)=\eta(1+t/\epsilon_L)
                  \eta(1+(t_{a3}-t)/\epsilon_R)
\end{equation}
leaves them exactly unchanged on $[0,t_{a3}]$ and makes them
smooth and compact in time as well as space. Flatness at the two
outer endpoints proves smoothness of their zero extensions.
Every additional mixed derivative follows from the ordinary
product rule, with the explicit factors
$\epsilon_L^{-b}[\eta^{(b)}]_0$ and
$\epsilon_R^{-b}[\eta^{(b)}]_0$.
The spatial fields solve the prescribed equations on the required
slab $[0,t_{a3}]$; no equation outside that slab is asserted after
the force cutoff. This proves the complete compact realization
through the actual third threshold and retains the precise finite
scope of the result.
